\begin{abstract}
Neural-graph programs intended for hardware require semantics that are deterministic and testable across multiple implementations and tool flows. NEMA defines a tick-based execution model for sparse neural graphs with an explicit snapshot rule, a fixed-point numeric model (RNE rounding and saturating overflow), and a bit-exact equivalence criterion based on per-tick SHA-256 digests of packed Q8.8 voltages. NEMA programs compile through a validated IR (mirrored by a proto schema) and can be exercised end-to-end via scripts that emit machine-readable reports.

We validate bit-exactness on two core benchmarks (B1 and B3) by running \texttt{nema bench verify}, which reports \texttt{ok=true} and \texttt{mismatches=[]} (mismatch=0) for both. We additionally exercise an external connectome-bundle workflow (B4) through the same manifest-driven interface. We provide hardware evidence through a self-audit gate: \texttt{audit\_min --mode hardware} returns \texttt{decision=GO}, indicating toolchain detection, HLS execution and QoR capture, and Vivado implementation timing with non-null WNS recorded in \texttt{bench\_report.json}. The artifact pack includes stable scripts, report paths, and a hashed manifest for reproducibility.
\end{abstract}
